\chapter{Đáp án tự kiểm tra}
\label{app:answers}

\section*{Chương 1 --- Spring Boot}
\begin{enumerate}
  \item Biến môi trường thắng: biến môi trường đứng trên tài liệu của
  config-server, và tài liệu đó đứng trên YAML đóng gói trong jar. Đúng
  với container vì image phải bất biến giữa các môi trường; những gì
  khác nhau do runtime bơm vào.
  \item \code{@Configuration}, \code{@ComponentScan},
  \code{@EnableAutoConfiguration}. Cái cuối cùng: auto-configuration
  phụ thuộc vào classpath, nên gỡ một dependency sẽ âm thầm gỡ luôn các
  bean mà nó kích hoạt.
  \item Nó được giải thành chuỗi ký tự \code{placeholder} theo đúng
  nghĩa đen. Vì thế service vẫn \emph{khởi động} được dù không có
  thông tin xác thực Google --- OAuth chỉ hỏng lúc người dùng bấm, chứ
  không hỏng lúc boot.
  \item Việc mở (exposure) khiến endpoint tồn tại; Spring Security vẫn
  phân quyền từng request. Nếu \code{/actuator/**} không được cho phép,
  probe nhận 401, và kubelet tính đó là thất bại.
\end{enumerate}

\section*{Chương 2 --- Config Server}
\begin{enumerate}
  \item Biến môi trường $>$ config-server $>$ YAML đóng gói.
  \item Native phục vụ file từ thư mục classpath bên trong jar; git
  thêm khả năng kiểm toán, review, hoàn tác, và thay đổi cấu hình mà
  không cần build lại config-server.
  \item \code{optional:} cho phép service boot khi config-server sập
  --- nhưng khi đó nó boot với các giá trị mặc định localhost và hỏng
  về sau, ở rất xa nguyên nhân.
  \item Trong \code{application.yml} của chính nó, đóng gói trong jar
  của nó. Nguyên tắc bootstrap: nguồn cấu hình không thể tự cấu hình
  cho mình; mọi chuỗi như vậy đều cần một gốc tự chứa.
\end{enumerate}

\section*{Chương 3 --- Eureka}
\begin{enumerate}
  \item course-service đăng ký (POST) lúc khởi động; discovery client
  của gateway làm mới cache registry của nó (mặc định tối đa
  30\,s); rồi việc định tuyến cân bằng tải từ cache. Trường hợp xấu
  nhất để một phiên bản mới nhận được lưu lượng: khoảng đăng ký + làm
  mới cache, tức hàng chục giây.
  \item Compose khiến \emph{tên} service phân giải được, và hostname
  trùng với tên đó; trong Kubernetes hostname của pod là tên pod, thứ
  không ai phân giải được --- nên phải đăng ký IP thay vào đó.
  \item Ưu tiên tính sẵn sàng của dữ liệu registry hơn tính chính xác.
  Có lợi: một switch khởi động lại làm im bặt mọi heartbeat nhưng các
  phiên bản vẫn đang chạy. Có hại: sập thật hàng loạt để lại các phiên
  bản đã chết vẫn nằm trong registry.
  \item Danh sách endpoint của Service thay cho registry; readiness
  probe thay cho heartbeat; kube-proxy/DNS thay cho tra cứu phía
  client; ReplicaSet thay cho việc phiên bản mới tự đăng ký.
\end{enumerate}

\section*{Chương 4 --- API Gateway}
\begin{enumerate}
  \item Trình duyệt $\rightarrow$ ts.local (file hosts) $\rightarrow$
  Traefik (Ingress: \code{/api} $\rightarrow$ Service api-gateway)
  $\rightarrow$ pod gateway (predicate \code{/api/auth/**},
  StripPrefix=2, không đòi JWT với \code{/api/auth}) $\rightarrow$
  pod auth-service do Eureka phân giải, đến nơi dưới dạng \code{/login}
  $\rightarrow$ controller.
  \item Google kiểm tra redirect URI từng byte một; nó đã được đăng
  ký với đường dẫn đầy đủ \code{/login/oauth2/...}, nên phải đến nơi
  nguyên vẹn. Các route API thì cắt tiền tố vì controller được viết
  không có tiền tố.
  \item Mọi request mang token do auth-service ký đều trượt bước kiểm
  tra của gateway --- 401 đồng loạt trên các route được bảo vệ trong
  khi bản thân việc đăng nhập có vẻ vẫn chạy.
  \item Proxy là bài toán I/O với độ đồng thời cao --- event loop rất
  hợp. course-service làm việc JDBC kiểu blocking; reactive sẽ làm nó
  phức tạp mà không có lợi ích gì cho đến khi chứng minh được điều
  ngược lại.
\end{enumerate}

\section*{Chương 5 --- Bảo mật}
\begin{enumerate}
  \item CSRF bám vào một thông tin xác thực có sẵn trong môi trường
  (cookie phiên được gửi tự động). JWT trong header Authorization được
  code gắn vào một cách tường minh; một request giả mạo xuyên site
  không thể thêm nó vào.
  \item Mười lăm phút (thời hạn của nó). Không gì cả --- chữ ký không
  thể thu hồi. \code{/refresh}: refresh token được đối chiếu với cơ sở
  dữ liệu, nơi việc thu hồi được lưu.
  \item Bất biến: mọi đường mạng đến một service đều đi qua filter JWT
  của gateway. Bị phá khi: một Ingress định tuyến thẳng đến service
  (route \code{/courses} đã xóa), hoặc bất kỳ bên gọi nào trong cụm
  giả mạo header trên mạng pod phẳng.
  \item Client-secret không bao giờ rời khỏi server; authorization
  code thì đi qua trình duyệt. Code chỉ dùng một lần, sống ngắn, và vô
  giá trị nếu không có client-secret lúc đổi lấy token.
  \item HS256: một secret dùng chung; ai kiểm tra được thì cũng ký
  được. RS256: khóa riêng ký, khóa công khai kiểm tra --- nên một IdP
  có thể cho hàng nghìn service kiểm tra token mà không service nào
  trong số đó đúc được token.
\end{enumerate}

\section*{Chương 6 --- Các service và module common}
\begin{enumerate}
  \item (a) HTTP --- bên gọi cần dữ liệu ngay. (b) Sự kiện --- một sự
  thật mà bên khác phản ứng theo; việc chấm điểm không được chờ mail.
  (c) Không cái nào: kiểm tra token là phép mật mã cục bộ với secret
  dùng chung; không cần gọi đi đâu.
  \item auth-service chỉ \emph{sản xuất} --- từng lần gửi hỏng riêng
  lẻ, boot vẫn thành công. notification-service chỉ \emph{tiêu thụ};
  không có broker thì nó chỉ là chi phí đứng không, nên nó đi kèm
  Kafka.
  \item Thuộc về common: các record sự kiện --- hai bên phải thống nhất
  hình dạng. Không bao giờ: logic nghiệp vụ/entity --- hành vi dùng
  chung trói chu kỳ phát hành của các service vào nhau và che mất ranh
  giới.
  \item Consumer triển khai trước producer, và thay đổi chỉ được thêm
  vào (trường mới, tùy chọn). Khi đó consumer cũ bỏ qua trường mới và
  consumer mới chịu được sự kiện cũ.
\end{enumerate}

\section*{Chương 7 --- Postgres, JPA, Flyway}
\begin{enumerate}
  \item Schema của hai lập trình viên âm thầm lệch nhau (update áp
  dụng theo từng cơ sở dữ liệu, không có lịch sử); và thay đổi phá hủy
  là bất khả thi --- update không bao giờ drop/đổi tên, nên refactor
  entity lặng lẽ lệch khỏi schema.
  \item Mỗi dòng lịch sử lưu checksum của script tương ứng; script bị
  sửa sẽ lệch checksum và Flyway từ chối khởi động --- lịch sử là dấu
  vết kiểm toán, và việc sửa sẽ làm giả nó.
  \item Phát hành 1: thêm cột mới, ghi cả hai (hoặc dùng trigger).
  Migration backfill. Phát hành 2: đọc cột mới, ngừng ghi cột cũ.
  Phát hành 3: drop cột cũ. Mỗi bước tương thích với phiên bản ngay
  trước nó.
  \item Ủng hộ: cô lập lỗi và phạm vi ảnh hưởng --- một server sập,
  nâng cấp, hay cấu hình sai không thể chạm tới dữ liệu của server
  kia. Phản đối: gấp đôi bộ nhớ và bề mặt vận hành với một homelab.
  Cloud có quản lý thường chọn cô lập bằng thông tin xác thực trên
  một instance vì server là đơn vị tính tiền.
\end{enumerate}

\section*{Chương 8 --- MongoDB}
\begin{enumerate}
  \item Mục từ điển: tự chứa, đọc/ghi như một khối, lồng nhau tự
  nhiên, không cần nhất quán giữa các mục. Enrollment: tham chiếu hai
  aggregate khác và đòi hỏi nhất quán giữa chúng --- join và
  transaction chính là điểm mấu chốt.
  \item Vào code ứng dụng (và validator của collection nếu muốn). Kỷ
  luật: lớp document đọc ngược được (trường mới là tùy chọn), và script
  sửa dữ liệu tường minh khi đổi hình dạng.
  \item Postgres từ chối kết nối đến cơ sở dữ liệu không tồn tại;
  Mongo tạo cơ sở dữ liệu ở lần ghi đầu tiên, nên tên sai cho ra kết
  quả rỗng. Ồn ào tốt hơn --- lỗi bạn thấy lúc triển khai hơn hẳn lỗi
  bạn thấy lúc kiểm toán.
  \item Chỉ ClusterIP + không có Ingress nghĩa là chỉ bên gọi trong cụm
  mới chạm tới được; sự tương đồng với compose cũng giữ ma sát khi dev
  bằng không. Mất hiệu lực ngay khi nó bị lộ ra ngoài cụm hoặc dùng
  chung cụm với workload không đáng tin.
\end{enumerate}

\section*{Chương 9 --- MinIO / S3}
\begin{enumerate}
  \item Service ký (micro giây) thay vì stream (giây); storage phục vụ
  byte trực tiếp. Kiểu lỗi mới: URL nhúng sẵn endpoint và thời hạn ---
  endpoint sai hay lệch đồng hồ làm hỏng tải xuống chỉ ở trình duyệt,
  service không nhìn thấy.
  \item Service ký chạm tới MinIO qua \code{minio:9000} (DNS của cụm);
  trình duyệt không phân giải được tên đó. Chữ ký bao trùm cả host,
  nên viết lại sau khi ký sẽ làm chữ ký vô hiệu --- vì thế phải ký
  với một endpoint thứ hai, công khai.
  \item Thời hạn (15 phút) giới hạn cửa sổ thiệt hại; URL chỉ cấp một
  object, không phải cả bucket; và không thể ký lại nếu không có khóa.
  \item Thay đổi: endpoint, thông tin xác thực (lý tưởng là chuyển
  sang IAM role), bảo đảm về độ bền. Không đổi: mọi lời gọi SDK, logic
  bucket/key/presign --- API chính là chuẩn.
\end{enumerate}

\section*{Chương 10 --- Kafka cốt lõi}
\begin{enumerate}
  \item Key theo userId --- thứ tự theo key chính là thứ tự theo
  partition. Số partition $\geq$ số consumer tối đa bạn từng muốn có
  trong group; độ song song bị chặn trên bởi số partition.
  \item Các group độc lập nhau: mỗi group giữ offset riêng cho từng
  partition, nên cả hai group đều duyệt hết toàn bộ log; tiêu thụ là
  đọc không phá hủy cộng với một con trỏ theo từng group.
  \item \code{LISTENERS} = socket được bind; \code{ADVERTISED\_LISTENERS}
  = địa chỉ \emph{trả về cho client} trong metadata. Bootstrap dùng
  địa chỉ bạn cung cấp; mọi kết nối sau đó dùng địa chỉ được quảng bá
  --- nên lỗi chỉ bắt đầu sau một cái bắt tay thành công.
  \item Retry và sập sau khi xử lý đều gây gửi lại; broker không thể
  biết bạn đã hành động chưa. Notification: khử trùng lặp theo id sự
  kiện (hoặc chấp nhận một mail trùng hiếm hoi --- lũy đẳng bằng cách
  chịu đựng).
\end{enumerate}

\section*{Chương 11 --- Spring Kafka}
\begin{enumerate}
  \item \code{send()} đệm ở phía client và trả về một future; không
  cần broker cho đến lúc giao. Sự cố lộ ra theo từng lần gửi: future
  thất bại sau timeout, và sự kiện đang đệm mất trắng nếu tiến trình
  chết.
  \item Khi chưa có offset đã commit, \code{earliest} bắt đầu từ record
  cũ nhất còn được giữ --- lịch sử được xử lý lại. Một khi group đã
  commit offset, thiết lập này không bao giờ được hỏi đến nữa.
  \item Không có: giải tuần tự ném lỗi ngay trong poll, offset không
  bao giờ tiến, partition kẹt mãi ở một record. Có: lỗi được vật chất
  hóa thành một thất bại có kiểu, error handler bỏ qua (hoặc retry rồi
  chuyển tiếp) --- ở production đầy đủ, chuyển đến
  \code{<topic>.DLT}.
  \item Giải tuần tự JSON khởi tạo các lớp được nêu tên trong header;
  danh sách tin cậy không giới hạn cho phép một message bắt consumer
  dựng lên kiểu bất kỳ trên classpath --- một vector tiêm gadget.
  \item Khi ``ghi DB + phát sự kiện'' không được phép lệch nhau. Nó
  làm cho thay đổi trạng thái và \emph{việc ghi lại ý định phát} trở
  thành nguyên tử --- một transaction cục bộ; một relay đảm bảo việc
  phát cuối cùng sẽ diễn ra.
\end{enumerate}

\section*{Chương 12 --- Docker}
\begin{enumerate}
  \item Namespace (góc nhìn cô lập) và cgroup (giới hạn tài nguyên).
  Không có kernel khách: khởi động tính bằng mili giây và không tốn chi
  phí hệ điều hành cho từng instance --- nhưng đồng thời: kernel dùng
  chung (cô lập yếu hơn VM) và JVM nhìn thấy tài nguyên của host trừ
  khi được bảo khác đi.
  \item Bên trong container, \code{localhost} là chính container đó.
  Compose sửa bằng DNS theo tên service trên mạng bridge; Kubernetes
  bằng DNS của Service --- cả hai qua biến môi trường ghi đè các giá
  trị mặc định localhost.
  \item Mỗi container có network namespace riêng, nên cả hai đều bind
  cổng 5432 ``của mình''. Xung đột quay lại ở host dùng chung: cổng
  host được publish (vì thế mới có 5433/5434) --- và với
  \code{--network host}, nơi không hề có sự tách biệt namespace.
  \item \code{depends\_on} thường chờ \emph{đã khởi động};
  \code{service\_healthy} chờ healthcheck vượt qua --- ở đây là
  config-server tự curl \code{/actuator/health} của chính nó, được mở
  bởi cấu hình management.
\end{enumerate}

\section*{Chương 13 --- Jib và Registry}
\begin{enumerate}
  \item Không cần daemon trong CI, chia layer theo độ biến động (push
  lại nhỏ), image không chạy bằng root và sẵn sàng cho OpenShift. Về
  bảo mật: không có daemon nghĩa là không có docker socket trong CI
  --- socket đó là root trên host.
  \item Tên image đồng thời là địa chỉ pull cho cụm. Ở đây chúng trùng
  nhau vì Jenkins dùng chung network namespace với host; nếu Jenkins
  chạy trên mạng bridge, \code{localhost:5000} của nó sẽ là chính nó
  và push hỏng trong khi pull vẫn chạy.
  \item Mỗi lần deploy phải import trên từng node, mất hàng phút tuần
  tự hóa, và không có lịch sử tag bất biến; lệnh \code{set image} của
  pipeline trỏ đến tên gắn tag SHA sẽ không có gì để trỏ vào ---
  registry chính là thứ làm cho tag có địa chỉ.
  \item Cùng tag, nội dung mới: với \code{IfNotPresent} image đã cache
  trên node âm thầm thắng. Cách sửa: \code{imagePullPolicy: Always}
  (lựa chọn của manifest), hoặc deploy bằng tag/digest bất biến (lựa
  chọn của pipeline) --- dự án này làm cả hai.
\end{enumerate}

\section*{Chương 14 --- Giải phẫu cụm}
\begin{enumerate}
  \item Bạn viết trạng thái mong muốn; controller điều hòa trạng thái
  thực tế về phía đó mãi mãi. Apply lần thứ hai không tạo ra diff, nên
  không có hành động --- lũy đẳng nhờ kiến trúc.
  \item Thứ tự là khái niệm thủ tục; hội tụ thay thế nó: mọi thứ cùng
  khởi động, bên phụ thuộc sập rồi thử lại với backoff cho đến khi thứ
  nó cần tồn tại. CrashLoop trong lúc hội tụ là cơ chế, không phải
  bug.
  \item kubectl gửi một lệnh REST delete đến API server; kho trạng
  thái cập nhật; vòng lặp của controller ReplicaSet thấy replica 0 $<$ 1
  và tạo một object pod; scheduler gán nó; kubelet trên node khởi động
  container.
  \item Apply hợp nhất trạng thái mong muốn; nó không bao giờ xóa
  những gì vắng mặt trong đầu vào. \code{kubectl delete ingress
  course-service -n ts}, hoặc apply với \code{--prune} (hoặc một
  controller GitOps) thì sẽ xóa.
\end{enumerate}

\section*{Chương 15 --- Workload}
\begin{enumerate}
  \item Mỗi thay đổi template sinh ra một ReplicaSet mới trong khi các
  ReplicaSet cũ được giữ ở scale 0 --- cập nhật cuốn chiếu và rollback
  một lệnh là hệ quả của việc tách bộ đếm khỏi template có phiên bản.
  \item Deployment: N replica thay thế được cho nhau, cái nào chết
  cũng được, thứ tự không quan trọng. StatefulSet: replica có tên với
  storage theo từng danh tính và thay thế có thứ tự. Eureka: Deployment
  (registry được dựng lại nhờ đăng ký lại). Redis cache: Deployment
  nếu thật sự chỉ là cache. Broker Kafka: StatefulSet. Gateway:
  Deployment.
  \item Pod mới phải vượt qua readiness trước khi pod cũ bị kết thúc
  --- probe của Chương~17 là cổng gác cho rollout của Chương~15.
  \item Cập nhật Deployment chạy pod cũ và mới đồng thời; hai tiến
  trình postgres cùng mở một thư mục dữ liệu sẽ làm hỏng nó.
  StatefulSet thay thế từng cái một và gắn lại claim cho pod kế nhiệm,
  không bao giờ hai cái cùng lúc.
\end{enumerate}

\section*{Chương 16 --- Cấu hình và Secret}
\begin{enumerate}
  \item RBAC về việc ai được đọc object Secret; tách vòng đời khỏi git;
  mã hóa khi lưu (tùy chọn). Dự án này: RBAC có (Role của Jenkins bỏ
  qua quyền đọc secret... nó cấp create/patch nhưng kiểm soát thực tế
  là giới hạn theo namespace), tách khỏi git có (script), mã hóa khi
  lưu không.
  \item Image server đọc \code{POSTGRES\_*}; Spring đọc
  \code{SPRING\_DATASOURCE\_*}. Tách Secret nghĩa là mỗi pod chỉ nhập
  những tên nó hiểu --- không có biến lạ che mất lỗi gõ hay rò rỉ
  phạm vi.
  \item Cập nhật: \code{JWT\_SECRET=new ./create-secrets.sh}; rồi ngay
  lập tức \code{kubectl -n ts rollout restart deployment/auth-service
  deployment/api-gateway}. Làm nửa chừng: bên ký và bên kiểm tra giữ
  khóa khác nhau --- mọi lần đăng nhập đều 401, hoặc tệ hơn, chỉ sau
  lần restart không liên quan tiếp theo.
  \item Cây \code{configurations/*.yml} của config-server, mount dưới
  dạng file --- và cho config-server nghỉ hưu luôn.
\end{enumerate}

\section*{Chương 17 --- Probe và tài nguyên}
\begin{enumerate}
  \item Startup thất bại quá ngân sách: restart. Readiness thất bại: bị
  gỡ khỏi endpoint, không restart --- có thể mạnh tay. Liveness thất
  bại: restart --- cái này phải dè dặt.
  \item \code{/v3/api-docs} chỉ trả lời sau khi toàn bộ context (kể cả
  Flyway) đã lên và được permitAll; một cổng đã bind chỉ chứng minh
  Tomcat tồn tại. \code{/actuator/health} trả 401 ở nơi security không
  cho phép (course/auth) và trung thực ở nơi cho phép (dictionary,
  gateway).
  \item OOM killer của kernel đã kết thúc cgroup vì vượt giới hạn bộ
  nhớ --- JVM không hề tham gia, nên không có exception. Kiểm tra:
  phép tính giữa limit và \code{MaxRAMPercentage}, và xem tải hay rò
  rỉ đã làm heap phình lên.
  \item Scheduler đọc requests (để xếp chỗ); kernel thực thi limits. Bộ
  nhớ không nén được --- vượt là có kẻ bị giết, nên hãy chặn nó một
  cách dự đoán được; CPU nén được --- limit sẽ bóp nghẹt lúc khởi động
  đúng khi JVM cần bứt tốc.
\end{enumerate}

\section*{Chương 18 --- Mạng}
\begin{enumerate}
  \item DNS phân giải tên Service thành một IP ảo ổn định; endpoint
  controller giữ danh sách pod Ready luôn cập nhật; kube-proxy (hoặc
  thứ tương đương của k3s) viết lại lưu lượng đến IP ảo sang IP của
  một pod được chọn; readiness quyết định tư cách thành viên.
  \item Chỉ pod Ready mới là endpoint, nên cập nhật cuốn chiếu chuyển
  lưu lượng đúng lúc pod mới vượt probe và pod cũ kết thúc --- không
  gián đoạn chính là tương tác giữa probe và endpoint.
  \item Selector/label lệch nhau --- Service khớp với không pod nào.
  Không có gì báo lỗi vì Service không có endpoint là một trạng thái
  hợp lệ; triệu chứng chỉ nhìn thấy trong \code{get endpoints}.
  \item \code{ts.local} không phải DNS thật; file hosts ghim nó vào IP
  của tailnet. Thay thế: một domain thật với DNS thực + TLS bằng
  cert-manager, hoặc một Cloudflare Tunnel công bố hostname đó.
\end{enumerate}

\section*{Chương 19 --- Lưu trữ}
\begin{enumerate}
  \item Pod mount PVC; PV (và StorageClass) có phạm vi toàn cụm; lớp
  gián tiếp này cho phép cùng một claim nhận local-path ở đây và EBS
  trên EKS --- manifest không đổi.
  \item Postgres chạy trên một node tại một thời điểm --- RWO là đủ.
  Upload dùng chung đẩy người ta sang RWX; lưu trữ đối tượng thường
  mới là câu trả lời đúng.
  \item Xóa StatefulSet không đụng đến PVC nên dữ liệu sống sót qua sai
  lầm điều phối; xóa PVC (reclaim \code{Delete}) gỡ luôn PV và dữ liệu
  --- dữ liệu sống ở claim, không phải ở workload.
  \item Bền vững (persistence): sống qua restart (có). Độ bền
  (durability): sống qua hỏng phần cứng (không --- một SSD). Sao lưu:
  sống qua sai lầm của chính bạn (không --- chưa có, khoảng trống
  thành thật).
\end{enumerate}

\section*{Chương 20 --- Kustomize}
\begin{enumerate}
  \item Một tập hợp chuẩn (không sót gì, một lệnh) và các chỉnh sửa
  xuyên suốt (namespace, label) được đóng dấu nhất quán.
  \item Namespace: phạm vi toàn cụm --- Role theo namespace của Jenkins
  không thể apply nó. Secret: giá trị tuyệt đối không được tồn tại
  trong git --- hình dạng nằm trong script, giá trị nằm trong môi
  trường của người vận hành.
  \item Kustomize: triển khai ứng dụng của chính bạn dưới dạng YAML dễ
  đọc. Helm: phân phối phần mềm có tham số đến cụm của người khác. Đến
  dưới dạng chart: Strimzi và kube-prometheus-stack.
  \item Không gì cả --- apply không dọn; mongodb vẫn chạy, giờ không
  còn ai quản lý. Đúng cách: \code{kubectl delete -f mongodb.yaml} (rồi
  gỡ khỏi danh sách), hoặc cơ chế prune/GitOps điều hòa cả việc xóa.
\end{enumerate}

\section*{Chương 21 --- Khái niệm CI/CD}
\begin{enumerate}
  \item CI: main luôn được xác nhận là build được --- stage Build. CD:
  mọi thay đổi đã xác nhận tự đến production không cần ai đỡ --- stage
  Deploy cộng với cổng gác rollout của nó.
  \item Jenkins nằm sau NAT, không có địa chỉ công khai --- GitHub không
  gọi vào được, nên Jenkins phải poll. Cùng ràng buộc: CI bên trong
  mạng doanh nghiệp mà git host trên cloud không chạm tới được.
  \item Build lại là một binary khác --- phân giải dependency, timestamp,
  toolchain khác --- nên test đã chứng nhận một thứ khác. Thứ tự ưu
  tiên biến môi trường của Chương~1 cho phép một artifact nhận cấu
  hình theo môi trường lúc chạy.
  \item Test bị bỏ qua (bật lại, hoặc gác ở chỗ khác và biết rõ điều
  đó); deploy không được xác minh (rollout status/smoke test); thành
  công được định nghĩa bằng exit code chứ không bằng hành vi (xác minh
  dựa trên probe, alert).
\end{enumerate}

\section*{Chương 22 --- Jenkins}
\begin{enumerate}
  \item Job, credential, trạng thái plugin, lịch sử build. Image =
  binary+plugin tái tạo được; volume = trạng thái không thay thế được;
  gộp lại, server là gia súc với một thư mục thú cưng.
  \item Jenkins chạy trên JDK\,21 của image gốc; build export
  \code{JAVA\_HOME=JDK25\_HOME} trong khối environment của Jenkinsfile
  --- theo từng pipeline, không phải toàn cục.
  \item apply -k cần get/list + create/patch/update trên mọi kind được
  apply; rollout status cần get/list/watch trên deployment và
  replicaset. Bỏ được nếu không có \code{set image}: không gì mới ---
  get+patch của set image đã được bao phủ sẵn; các quy tắc đọc pods/log
  chỉ tồn tại cho trình xử lý thất bại.
  \item Jib push đến \code{localhost:5000}, và tên đó phải chỉ cùng
  một máy với cả bên push lẫn bên pull; host networking khiến
  localhost của Jenkins chính là localhost của host, đúng như tên
  image đã hứa với k3s.
  \item ServiceAccount = danh tính; Role = danh sách năng lực;
  RoleBinding = mối nối. Role chứ không phải ClusterRole giới hạn một
  pipeline bị xâm nhập trong một namespace --- phạm vi ảnh hưởng như
  một lựa chọn về kiểu.
\end{enumerate}

\section*{Chương 23 --- Pipeline}
\begin{enumerate}
  \item SNAPSHOT là khả biến --- một node có thể phục vụ bất kỳ build
  lịch sử nào dưới cùng một tên; SHA ghim đúng commit và làm cho câu
  hỏi ``cái gì đang chạy?'' trả lời được từ \code{kubectl describe}.
  \item Chỉ build các module được liệt kê cộng với dependency của
  chúng. Quên một mục MODULES: service mới không bao giờ có image;
  vòng lặp deploy sau đó thất bại (hoặc deploy một image cũ nếu có)
  --- điểm ghép nối thủ công duy nhất của pipeline.
  \item \code{kubectl rollout status --timeout=300s} --- nó chuyển kết
  quả probe (Chương~17) thành exit code của pipeline.
  \item Mọi rollout nên tiến hành đồng thời; chờ theo từng service
  tuần tự hóa sự hội tụ mà Kubernetes sẵn sàng làm song song --- cùng
  tổng công việc, thời gian thực dài hơn.
  \item Thay đổi: cách kích hoạt (webhook), nơi build chạy (runner tạm
  thời), cơ chế secret (secret của repo). Không đổi: thứ tự stage, lệnh
  Maven, gắn tag SHA, bước xác minh rollout.
\end{enumerate}

\section*{Chương 24 --- Vận hành}
\begin{enumerate}
  \item get pods $\rightarrow$ describe pod $\rightarrow$ logs
  (\code{--previous} --- thiết yếu với CrashLoop, khi container hiện
  tại chưa có lịch sử gì) $\rightarrow$ exec vào và thăm dò từ bên
  trong.
  \item CrashLoop: tiến trình thoát --- tầng ứng dụng.
  \code{Running 0/1}: tiến trình sống nhưng readiness thất bại ---
  đích của probe (thường là một dependency) mới là tầng cần xem.
  \item Exec vào pod và thử kết nối từ đó --- DNS, đường mạng, và biến
  môi trường được bơm vào của pod khác với laptop của bạn; laptop chỉ
  chứng minh Postgres còn sống, không chứng minh pod chạm tới được nó.
  \item Lưu lượng pod/service của k3s đi qua firewall của host; chặn
  \code{10.42.0.0/16} hoặc \code{10.43.0.0/16} làm hỏng DNS và các
  lời gọi pod-đến-pod một cách chập chờn --- triệu chứng không phân
  biệt được với ứng dụng thiếu ổn định.
\end{enumerate}

\section*{Chương 25 --- Kafka trong cụm}
\begin{enumerate}
  \item Một controller có vòng lặp điều hòa nhắm đến các kind tài
  nguyên tùy chỉnh do CRD định nghĩa. Strimzi thêm
  \code{Kafka}/\code{KafkaNodePool} (cụm) và \code{KafkaTopic} (các hợp
  đồng) --- cộng thêm các kind user, connect và bridge mà dự án này
  không dùng.
  \item Gói operator chứa các object phạm vi toàn cụm (CRD, ClusterRole)
  mà Role theo namespace không thể cấp; \code{kafka.yaml} chỉ chứa tài
  nguyên tùy chỉnh theo namespace, thứ Role giờ đã cho phép. Role hay
  ClusterRole chính là ranh giới.
  \item Advertised listener, các voter của controller quorum, bản đồ
  giao thức bảo mật theo từng listener (thêm cả cluster ID và Service
  cho từng broker). Operator tính chúng từ khai báo listener và node
  pool.
  \item Trạng thái thực của pod broker (\code{get pods}) và các dòng
  điều hòa trong log của operator. Condition có thể là một lần ghi cũ
  từ một vòng điều hòa đã timeout trong lúc pull image kéo dài; vòng
  lặp định kỳ tiếp theo sẽ sửa lại.
  \item Số partition được review, không có cuộc đua producer/topic ngày
  đầu, và một bảng kê các hợp đồng sự kiện dưới dạng object. Một
  partition: thứ tự toàn phần theo topic, đủ dùng chừng nào mỗi consumer
  group chỉ có một instance; partition có thể thêm về sau nhưng không
  bao giờ bớt được.
\end{enumerate}

\section*{Chương 26 --- Frontend và host đơn}
\begin{enumerate}
  \item Giai đoạn một (Node) biên dịch và đóng gói; giai đoạn hai (nginx)
  phục vụ. Chỉ \code{dist/} được chép bằng \code{COPY --from=build};
  Node, \code{node\_modules} và trình biên dịch bị bỏ đi, để lại một
  image $\sim$15\,MB với bề mặt tấn công nhỏ hơn.
  \item Traefik so \code{/courses/7} với cả hai Ingress; \code{/} là
  tiền tố khớp dài nhất (không có quy tắc \code{/courses} nào), nên
  nginx nhận nó. Không có file như vậy;
  \code{try\_files \$uri \$uri/ /index.html} rơi về
  \code{index.html}, ứng dụng khởi động và React Router hiển thị route
  từ URL.
  \item Các quy tắc cho cùng một host được hợp nhất vào một bảng định
  tuyến và tiền tố đường dẫn khớp dài nhất thắng: \code{/api/courses}
  khớp \code{/api} (dài hơn) thay vì \code{/}; \code{/assets/x.js} chỉ
  khớp \code{/}.
  \item CORS chỉ áp dụng cho request xuyên origin. Trong cụm, trang và
  API dùng chung scheme, host và cổng (\code{http://ts.local}), nên
  trình duyệt không bao giờ gửi preflight; trên laptop trang là
  \code{localhost:3000} còn API là \code{localhost:8080} --- cổng khác
  là origin khác.
  \item Socket thuộc sở hữu của group \code{docker} trên host (GID 984
  trên ts-server); user jenkins bên trong container không thuộc group
  nào có ID số đó, nên bị từ chối ghi. \code{--group-add} gắn GID của
  host vào lúc chạy; tính toán nó giữ cho script đúng trên mọi host mà
  không cần cứng hóa một con số đặc thù của máy.
\end{enumerate}

\section*{Chương 27 --- Chặng đường phía trước}
\begin{enumerate}
  \item Độ trễ tiêu thụ (consumer-group lag): khoảng cách giữa offset
  mới nhất trong một partition và offset group đã commit --- lag tăng
  nghĩa là notification-service đang tụt lại hoặc đã sập.
  \item Đẩy: thông tin xác thực nằm ở CI (kubeconfig của Jenkins). Kéo:
  nó không bao giờ rời khỏi cụm; CI chỉ ghi vào git. Prune giải quyết
  cái bẫy apply-không-bao-giờ-xóa của Chương~20.
  \item Route của gateway đổi \code{lb://name} thành
  \code{http://name:port}; endpoint của Service thay cho đăng ký,
  readiness probe thay cho heartbeat; xóa discovery-server (và các
  dependency eureka client), giải phóng tài nguyên của nó.
  \item Chúng đã chạy không bằng root với quyền của group 0, nên SCC
  UID-tùy-ý của OpenShift chấp nhận chúng nguyên vẹn; chỉ còn
  Ingress $\rightarrow$ Route.
\end{enumerate}

\section*{Chương 28 --- Đào sâu: JWT ở quy mô lớn}
\begin{enumerate}
  \item Không. Config-server giờ chỉ giữ khóa công khai, thứ xác minh
  được nhưng không ký được. Giả mạo cần PEM riêng, chỉ nằm trong Secret
  của auth-service --- Secret đó (hoặc pod auth-service) mới là thứ
  phải lộ.
  \item Công bố khóa công khai mới trong JWKS dưới một \code{kid} mới
  trong khi giữ khóa cũ; đợi cache của các bên xác minh làm mới; chuyển
  ký sang khóa riêng mới; sau thời gian sống dài nhất của token, gỡ
  khóa cũ. Ở bước~3 hai thế hệ token lưu hành cùng lúc; \code{kid} cho
  mỗi bên xác minh biết \emph{khóa công khai nào} cần áp dụng, nên
  token cũ vẫn xác minh được và khóa cũ về sau có thể nghỉ hưu một cách
  chắc chắn chứ không phải bằng cách thử.
  \item User đã cầm access token không bị ảnh hưởng cho tới khi nó hết
  hạn (15 phút), vì các bên xác minh đã cache khóa công khai. Sau đó
  \code{/refresh} thất bại và họ bị đăng xuất; đăng nhập mới thất bại
  ngay từ phút đầu. Một bên xác minh khởi động lại trong lúc sự cố boot
  với cache JWKS rỗng và, khi fail closed, từ chối mọi token --- đó là
  lý do khóa công khai cũng nên có sẵn trong cấu hình làm phương án dự
  phòng.
  \item Lỗi thời: claim bị đóng băng lúc phát hành, nên một khóa học
  tạo sau khi đăng nhập là vô hình cho tới khi token hết hạn --- token
  không thể biết dữ liệu nó không sở hữu. Quyền sở hữu: thành viên khóa
  học là dữ liệu của course-service; mã hóa nó vào token do auth-service
  phát hành là ghép schema của hai service qua hợp đồng token. (Kích
  thước là lý do thứ ba, và ít quan trọng nhất.)
  \item Header được chuyển tiếp, nhưng không endpoint công khai nào đọc
  \code{X-User-Id}: \code{/refresh} xác thực bằng refresh token,
  \code{/actuator} không xác thực gì. Nên hôm nay không khai thác được,
  nhưng bất biến ``chỉ gateway viết các header này'' không được cưỡng
  chế. \code{default-filters: RemoveRequestHeader=X-User-Id} (và
  \code{X-User-Role}) trong \code{api-gateway.yml} cưỡng chế nó.
  \item Tối đa 15 phút --- thời gian sống của access token --- vì
  gateway xác minh cục bộ và chỉ \code{/refresh} mới hỏi cơ sở dữ liệu.
  Deny-list theo \code{jti} thêm Redis (một lượt tra mỗi request, Redis
  chết thì fail-open hay fail-closed tùy bạn chọn); mẫu phantom token
  thêm Redis trên đường găng của mọi request mà không có lựa chọn
  fail-open --- Redis chết là sập toàn bộ.
\end{enumerate}

\section*{Chương 29 --- Đào sâu: Cấu hình}
\begin{enumerate}
  \item Secret thắng: nó tới dưới dạng biến môi trường
  \code{SPRING\_DATASOURCE\_PASSWORD}, và môi trường OS đứng trên tài
  liệu từ config-server. Nếu object Secret không tồn tại,
  \code{envFrom} làm pod thất bại; nếu chỉ thiếu khóa, pod khởi động,
  giá trị \code{root} nguyên văn trong file được phục vụ có hiệu lực,
  và service kết nối bằng \code{root}, lặng lẽ.
  \item Ưu tiên chọn \emph{tài liệu nào} định nghĩa
  \code{app.s3.secret-key}; tài liệu của config-server định nghĩa, và
  nó gửi nguyên văn \code{\${APP\_S3\_SECRET\_KEY:minioadmin}} chưa
  phân giải. Client phân giải placeholder với Environment của chính nó
  sau đó, nên biến môi trường thắng vì văn bản yêu cầu thế, không phải
  vì thứ tự nguồn.
  \item Cửa sổ một: giữa hai lần refresh, auth-service ký bằng khóa mới
  còn gateway xác minh bằng khóa cũ --- mọi lần đăng nhập cho ra token
  gateway từ chối. Cửa sổ hai: sau khi cả hai refresh, mọi token phát
  trước thay đổi bị từ chối cho tới lúc đằng nào cũng hết hạn. Xoay
  nóng cần các khóa chồng lấn chọn bằng \code{kid}, thứ một secret
  HS256 dùng chung không diễn đạt được.
  \item Không optional: lấy thất bại, JVM thoát, kubelet khởi động lại
  với back-off tăng dần cho tới khi server trở lại. \code{optional:}:
  boot từ mặc định đóng gói trỏ về localhost, qua readiness probe, nhận
  lưu lượng, thất bại mọi request. \code{fail-fast} cộng retry: thử lại
  trong tiến trình với back-off trong cửa sổ đã cấu hình, rồi thoát ---
  không lãng phí lần khởi động JVM nào, và không có pod nói dối.
  \item Generator của Kustomize gắn hash nội dung vào tên ConfigMap;
  file đổi là tên mới, Deployment tham chiếu nó là pod template mới, và
  Kubernetes roll các pod. Cấu hình, code và rollback thành một cơ chế
  có lịch sử --- server không cho được điều đó nếu không có Bus hay
  khởi động lại.
  \item Bỏ khi mọi runtime là Kubernetes, không giá trị nào phải đổi mà
  không khởi động lại pod, và secret đã sống trong nền tảng. Giữ khi
  một runtime ngoài Kubernetes (VM, Compose ở production) phải dùng
  chung cùng nguồn cấu hình.
\end{enumerate}

\section*{Chương 30 --- Đào sâu: Khám phá dịch vụ}
\begin{enumerate}
  \item Đăng ký tức thì khi web server khởi động (Spring ép việc đó).
  Rồi cache phản hồi chỉ đọc của server (30\,s), chu kỳ tải registry
  của gateway (30\,s), và TTL cache của LoadBalancer (35\,s): khoảng
  95\,s xấu nhất, thường bằng một nửa.
  \item Lease hết hạn 90\,s, rồi lượt quét loại bỏ (tới 60\,s), rồi
  cache server 30\,s, tải của client 30\,s, cache LoadBalancer 35\,s
  --- tới 245\,s. Tắt êm gửi DELETE và bỏ qua hai bước đầu (150\,s),
  chỉ còn lại các cache.
  \item Gia hạn kỳ vọng mỗi phút: $5 \times 2 = 10$; ngưỡng
  $\lfloor 10 \times 0{,}85 \rfloor = 8$; server chỉ khỏe khi gia hạn
  nhiều hơn 8. Một client chết cho đúng 8, nên tự bảo toàn kích hoạt
  và không gì bị loại bỏ. \code{eureka.server.enable-self-preservation:
  false} (hoặc \code{renewal-percent-threshold} thấp hơn) sửa việc đó
  trên đội tàu nhỏ.
  \item Trạng thái tự báo của Eureka gác đường của gateway, vì
  \code{lb://} phân giải thành IP pod đã đăng ký và không bao giờ chạm
  tới Service Kubernetes nơi readiness được cưỡng chế. Định tuyến tới
  \code{http://course-service:8082} đưa kube-proxy, và do đó readiness,
  trở lại đường đi.
  \item Compose không kích hoạt profile nào, nên tài liệu gốc áp dụng:
  Eureka client bật, route \code{lb://}, và \code{url} của Feign phân
  giải thành chuỗi rỗng, thứ Feign hiểu là ``dùng phân giải cân bằng
  tải theo tên''. Dưới \code{k8s} thuộc tính đó là URL Service và Feign
  bỏ qua LoadBalancer hoàn toàn.
  \item Registry biết một instance chỉ qua điều instance tự nói về
  mình, và biết đến sự im lặng chỉ sau một lease nó không quét ngay;
  endpoint của Service do kubelet duy trì, thứ chủ động probe pod vài
  giây một lần và gỡ nó ngay khoảnh khắc probe thất bại. Tự báo với
  thời gian ân hạn so với quan sát bên ngoài không có ân hạn.
\end{enumerate}

\section*{Chương 31 --- Đào sâu: Gateway dưới tải}
\begin{enumerate}
  \item Chưa có gì được gửi đi: \code{chain.filter} trả về một
  \code{Mono} mô tả phần còn lại của chuỗi; nó chỉ chạy khi framework
  subscribe, và thread gọi đã rảnh. \code{doFilter} của filter servlet
  chặn tới khi lời gọi phía sau hoàn tất, nên thread bị giữ suốt
  request --- chính mô hình mà event loop tồn tại để tránh.
  \item Mọi request có socket thuộc vòng lặp đang khựng đều dừng ---
  kể cả các route không hề chạm Redis --- vì vòng lặp dùng chung giữa
  các route, không phải theo route. Reactor Netty chạy
  \code{max(cores, 4)} vòng lặp; trên máy 2 lõi là bốn, nên một vòng
  lặp bị chặn là một phần tư năng lực gateway.
  \item Mười sáu: pool mặc định theo host là \code{max(cores, 8)
  $\times$ 2}. Không có timeout phản hồi, 16 request treo giữ cả 16 kết
  nối và request thứ 17 đợi 45\,s trong hàng acquire. Đặt
  \code{spring.cloud.gateway.httpclient.response-timeout} (hoặc tăng
  \code{pool.max-connections}) đổi con số; timeout mới là thứ sửa
  nguyên nhân.
  \item Filter bọc mỗi lời gọi trong một TimeLimiter Resilience4j có
  mặc định 1\,s. Lần tải lên 1,5\,s bị cắt, tính là thất bại, và sau đủ
  số lần, breaker mở chống lại một service khỏe mạnh. Đặt
  \code{timeout-duration} của TimeLimiter instance đó ít nhất dài
  bằng timeout phản hồi.
  \item Với hai replica trở lên, bucket trong bộ nhớ cho mỗi replica
  giới hạn riêng, nên một user được $N\times$ hạn ngạch. Redis giữ một
  bucket cho mỗi khóa, cập nhật nguyên tử bằng script Lua. Đánh đổi:
  Redis là phụ thuộc trên mọi request; mặc định của filter khi Redis
  không với tới là cho request qua (fail-open), việc phải là lựa chọn
  có chủ ý.
  \item Mọi request đều không đơn giản, nên khi dev mỗi lời gọi API đi
  sau một preflight \code{OPTIONS}. \code{setMaxAge} cho trình duyệt
  cache kết quả preflight (mặc định của Chrome là 5\,s). Trong cụm,
  frontend và \code{/api} chung một origin phía sau Traefik, nên không
  request nào xuyên origin và filter CORS không bao giờ được hỏi tới.
\end{enumerate}

\section*{Chương 32 --- Đào sâu: Gọi đồng bộ}
\begin{enumerate}
  \item (a) Không code nào gọi \code{AuthServiceClient} --- nối nó vào
  một service như \code{UserLookupService}. (b) Tích hợp breaker của
  Feign đang tắt --- đặt \code{circuitbreaker.enabled: true} dưới
  \code{spring.cloud.openfeign}. (c) Không
  có bean \code{CircuitBreakerFactory} --- thêm dependency
  \code{spring-cloud-starter-}\allowbreak\code{circuitbreaker-resilience4j}.
  Riêng nữa, không gì tham chiếu instance \code{authService} --- gắn
  \code{@CircuitBreaker} với \code{name = "authService"} lên phương
  thức gọi.
  \item Chín: \code{minimum-number-of-calls} mặc định bằng cỡ cửa sổ,
  nên breaker không đánh giá cho tới kết quả thứ mười. Lời gọi 30\,s là
  thành công vì \code{slow-call-duration-threshold} mặc định 60\,s;
  chỉ lời gọi chậm hơn thế mới tính là chậm, và chỉ khi đạt ngưỡng tỷ
  lệ chậm (mặc định 100\,\%).
  \item Trình duyệt (dài nhất) $>$ timeout phản hồi của gateway $>$
  Feign read timeout $\times$ số lần thử (cộng backoff). Mỗi chặng phải
  bỏ cuộc trước bên gọi của nó, nếu không nó làm việc cho request đã bị
  bỏ rơi; và số lần retry nhân với timeout trong cùng, nên phải nằm
  trong ngân sách của gateway --- 2\,s $\times$ 2 lần dưới timeout
  gateway 5\,s, không phải 3.
  \item Một \code{NullPointerException} ở mọi bên gọi quên kiểm tra,
  chỉ lộ ra khi có sự cố; và không phân biệt được 404 (user thực sự
  không có) với một lần breaker từ chối, nên cả hai bị đối xử như nhau.
  \code{degraded(Long id, Throwable cause)} --- fallback hai tham số
  --- nhận nguyên nhân và có thể trả giá trị đã cache hay giá trị thay
  thế cho lỗi hạ tầng trong khi để 4xx lan tiếp.
  \item \code{SecurityConfig} của auth-service chỉ cho phép đường dẫn
  công khai; \code{/users/\{id\}} cần ngữ cảnh đã xác thực mà
  \code{GatewayHeaderAuthFilter} của nó dựng từ \code{X-User-Id}, thứ
  Feign không gửi. Mỗi 401 là một ngoại lệ, nên mười lần mở breaker
  chống lại một service khỏe. Sửa: một \code{RequestInterceptor} của
  Feign chuyển tiếp header danh tính, và \code{ignore-exceptions} cho
  4xx để lỗi phía client không bao giờ tính là lỗi phụ thuộc.
  \item $0{,}999^3 \approx 99{,}7\,\%$. Nâng lên bằng cách bỏ phụ
  thuộc khỏi đường request --- bản sao cục bộ nuôi bằng sự kiện, hoặc
  cache đặt trước lời gọi --- và làm các phụ thuộc còn lại thành tùy
  chọn với fallback có kiểu để lỗi của chúng làm suy giảm response thay
  vì làm nó thất bại.
\end{enumerate}

\section*{Chương 33 --- Đào sâu: Dữ liệu xuyên service}
\begin{enumerate}
  \item Thứ tự A: transaction commit, rồi JVM chết trước khi bộ đệm
  producer flush --- user đã kích hoạt, không có sự kiện, không bao giờ
  có học viên. Thứ tự B: bộ đệm flush và broker ack, rồi commit thất
  bại --- course-service tạo học viên cho một user đã bị rollback. B
  tạo ra dòng ma.
  \item Producer không lấy được metadata, chặn trong
  \code{max.block.ms} (60\,000\,ms), rồi ném ngoại lệ; ngoại lệ thoát
  khỏi \code{register()}, transaction rollback, và user không được tạo.
  Với outbox, request ghi một dòng user và một dòng outbox trong một
  commit và trả về sau vài mili giây; relay thử lại việc gửi ở nền.
  \item Hai replica chạy cùng scheduler. Không có khóa, cả hai chọn
  cùng các dòng chưa phát và cả hai phát chúng --- một bản sao cho mọi
  sự kiện, mỗi giây. \code{FOR UPDATE} bắt replica thứ hai chờ;
  \code{SKIP LOCKED} bắt nó lấy các dòng kế tiếp thay vào đó, nên chúng
  chia việc mà không chồng lấn.
  \item Thứ tự theo key. Nếu dòng~1 của user~42 thất bại và vòng lặp
  \code{continue} sang dòng~2 cùng user, dòng~2 được phát trước; nhịp
  kế phát dòng~1 phía sau nó, và consumer thấy ``updated'' trước
  ``activated''. Dừng lô giữ thứ tự partition bằng thứ tự outbox.
  \item Không gì: \code{ts.user.activated} bắn một lần, và không có
  \code{ts.user.updated}, nên \code{student.email} đóng băng từ lúc
  kích hoạt. Thay đổi: phát một dòng outbox với
  \code{eventType=UserUpdated} ở bất cứ đâu lưu hồ sơ; làm
  \code{UserActivatedListener} upsert (tìm-hoặc-tạo, rồi đặt trường);
  cho nó đăng ký cả hai topic hoặc thêm listener thứ hai.
  \item Trạng thái: \code{ACTIVE} $\rightarrow$ \code{DELETING}
  $\rightarrow$ đã xóa, hoặc về lại \code{ACTIVE}. Sự kiện: auth phát
  \code{UserDeletionRequested} qua outbox; course phát
  \code{StudentDataDeleted} hoặc, khi hỏng,
  \code{StudentDataDeletionFailed} sau khi xóa học viên, ghi danh, bài
  nộp và object MinIO. Bù: khi
  thất bại auth khôi phục \code{ACTIVE} và cảnh báo. \code{DELETING} là
  khóa ngữ nghĩa: nó từ chối đăng nhập và ghi danh mới trong lúc saga
  chạy, nên không bước nào thao tác trên dữ liệu bước khác đang xóa.
\end{enumerate}

\section*{Chương 34 --- Đào sâu: Bảo đảm giao nhận Kafka}
\begin{enumerate}
  \item \code{enable.idempotence=true} (broker khử trùng các lần thử
  lại của producer), \code{acks=all}, \code{retries} không giới hạn
  trong \code{delivery.timeout.ms} --- toàn bộ là mặc định của client
  Kafka~3.x mà dự án chưa từng đặt. Trên Compose có một broker và
  \code{replicas: 1}, nên ``mọi replica đồng bộ'' là một cái đĩa; thiết
  lập này chỉ mua được độ bền khi \code{replication.factor=3} và
  \code{min.insync.replicas=2} tồn tại.
  \item Một lô được poll; listener gửi mail từng cái trên luồng poll;
  xử lý vượt \code{max.poll.interval.ms} (300\,000\,ms); broker đuổi
  consumer và gán lại partition từ offset đã commit cuối; chủ mới xử lý
  cùng các record trong khi chủ cũ vẫn chạy nốt vòng lặp.
  \item Nếu lượt giành commit một mình rồi lần ghi nghiệp vụ thất bại,
  lần thử lại thấy id đã được giành, trả về, và offset tiến lên --- thay
  đổi mất vĩnh viễn trong khi bảng nói nó đã xảy ra. Trong một
  transaction, thất bại rollback cả lượt giành, nên lần thử lại xử lý
  được.
  \item Chờ đợi chỉ giúp lỗi tạm thời: một SMTP relay quay lại, một
  nháy mạng. \code{NullPointerException} hay lỗi giải tuần tự là tất
  định; thử lại ba lần chỉ trì hoãn cùng thất bại. Nó bị loại trừ, nên
  đi thẳng tới DLT ngay lần thử đầu.
  \item \code{KafkaAutoConfiguration} của Boot lùi lại khi có bean
  \code{ProducerFactory}, nên \code{spring.kafka.producer.*} trong
  \code{course-service.yml} bị course-service bỏ qua. Sửa: xóa
  \code{KafkaProducerConfig} (factory của Boot đọc cùng các serializer
  từ YAML) --- hoặc giữ bean và xóa khối YAML.
  \item Retention là 168 giờ; chín ngày dài hơn. Offset đã commit của
  group giờ trỏ trước record cũ nhất còn giữ, nên
  \code{auto-offset-reset: earliest} có hiệu lực và nó tiếp tục từ record
  còn sống cũ nhất --- hai ngày sự kiện đã mất. Với group id mới không có
  offset commit nào; \code{earliest} khiến nó xử lý toàn bộ bảy ngày còn
  giữ, gửi lại mọi mail trong đó.
\end{enumerate}

\section*{Chương 35 --- Đào sâu: File ở quy mô lớn}
\begin{enumerate}
  \item \code{endpointOverride}: client dùng \code{endpoint}
  (\code{minio:9000}), presigner dùng \code{public-endpoint}. Nếu
  presigner ký với tên nội bộ, mọi link tải xuống sẽ nhúng một host trình
  duyệt không phân giải được --- và viết lại sau đó làm hỏng chữ ký.
  \item SigV4 ký một canonical request bao gồm các header được ký, và
  \code{Host} luôn nằm trong số đó (\code{X-Amz-SignedHeaders=host}).
  Link ký cho \code{minio:9000} rồi viết lại sang host khác trình ra
  một canonical request khác; MinIO tính lại chữ ký, không khớp, và trả
  về \code{403 SignatureDoesNotMatch}.
  \item Transaction chỉ bao phủ insert vào Postgres; \code{putObject}
  là một lời gọi HTTP bên ngoài nó. \code{save} thất bại thì dòng bị
  rollback trong khi object vẫn nằm trong bucket dưới một key không
  dòng nào ghi lại --- một object mồ côi.
  \item Rẻ: liệt kê mọi lần nộp của một học sinh cho một bài tập (một
  \code{ListObjectsV2} với tiền tố). Đắt: tìm mọi file của một học sinh
  qua các bài tập (không tiền tố nào khớp; quét toàn bộ hoặc truy vấn
  cơ sở dữ liệu). UUID khiến mọi key duy nhất, nên không upload nào ghi
  đè upload khác; versioning chỉ bảo vệ trước những lần ghi đè không
  thể xảy ra.
  \item Nộp text thành công (không gọi kho). Liệt kê thành công và trả
  về link (ký là tính toán cục bộ). Mở link thất bại (MinIO phải phục
  vụ nó). Khởi động lại course-service thất bại: kiểm tra bucket trong
  \code{@PostConstruct} ném lỗi khi connection refused.
  \item Key được dựng phía server từ \code{assignmentId} và
  \code{studentId} của chính người gọi, và lưu trên một dòng
  \code{PENDING\_FILE} thuộc học sinh đó; \code{completeUpload} kiểm tra
  dòng thuộc về người gọi và \code{HEAD} trên key của \emph{chính} dòng
  đó thành công. URL cho key của học sinh khác không gắn được vào dòng
  khác vì dòng không bao giờ nhận key do client cung cấp.
\end{enumerate}

\section*{Chương 36 --- Đào sâu: Từ điển trên MongoDB}
\begin{enumerate}
  \item \code{filter} của stage \code{FETCH} loại chúng. Index
  \code{createdByUserId} chỉ thu hẹp về lát của người dùng; một
  \code{\$regex} không neo, không phân biệt hoa thường không có ký tự
  đầu nào để seek, nên mọi document đã nạp đều bị thử. Neo pattern
  (\code{\^{}clou}) với collation phù hợp cho phép index trên
  \code{word} phục vụ nó như một khoảng.
  \item Mọi kết quả khớp --- có thể là cả 3{,}000 --- trở về từ Mongo;
  \code{toResponse} rồi chạy hai truy vấn \code{word\_links} mỗi kết
  quả (khoảng 6{,}000); trang được dựng trong Java bằng \code{subList}
  sau khi sort toàn bộ danh sách.
  \item Transaction nhiều document cần replica set; dự án chạy một
  \code{mongod} standalone. Crash sau khi xóa định nghĩa và trước khi
  xóa các liên kết để lại cạnh trỏ vào một từ đã mất.
  \code{findByIdAndCreatedByUserId(...).orElse(null)} rồi
  \code{filter(Objects::nonNull)} âm thầm bỏ các node như vậy.
  \item \code{BsonMaximumSizeExceededException} (giới hạn 16\,MB). Sáu
  ảnh khoảng 2--3\,MB mỗi cái, phình thêm một phần ba khi base64, vượt
  16{,}793{,}600 byte. Quy tắc của Chương~35: byte trong kho đối tượng,
  key trong cơ sở dữ liệu --- lưu ảnh trong MinIO và giữ key của chúng
  trong \code{images}.
  \item Text index tách token và rút gốc các từ trọn vẹn; ``lou'' không
  phải token của ``cloud'', nên không có kết quả. Khớp giữa từ cần quét
  regex hoặc index n-gram/trigram (Atlas Search, hay \code{pg\_trgm}
  trong Postgres).
  \item dictionary-service hỏng trên mọi endpoint; mọi service và topic
  khác không bị ảnh hưởng (không kết nối chung, không bên gọi, không sự
  kiện). Thêm một lệnh ping Mongo vào readiness probe sẽ đưa pod ra
  khỏi danh sách endpoint của gateway, biến 500 thành một 503 sạch sẽ ở
  biên.
\end{enumerate}

\section*{Chương 37 --- Đào sâu: Khả năng quan sát}
\begin{enumerate}
  \item Phía producer:
  \code{spring.kafka.\allowbreak template.\allowbreak observation-enabled}
  (vắng mặt trong mọi tài liệu cấu hình); phía consumer:
  \code{spring.kafka.\allowbreak listener.\allowbreak observation-enabled}.
  Thuộc tính chỉ cấu
  hình \code{KafkaTemplate} \emph{tự cấu hình}; auth-service dùng cái
  đó, nên thuộc tính sửa được nó. Course-service tự dựng template trong
  \code{KafkaProducerConfig}, thứ bỏ qua thuộc tính và cần
  \code{setObservationEnabled(true)} trong code.
  \item Nó tiếp tục trace: mặc định của Boot~3 là \emph{đọc} W3C, B3
  và B3-multi và \emph{phát} W3C, quyết định bởi các mặc định của
  \code{management.tracing.propagation}. Header rời gateway là
  \code{traceparent}, cùng trace id và một span id cha mới.
  \item Exposure là danh sách tên; endpoint chỉ tồn tại khi một
  registry hiển thị định dạng đó nằm trên classpath. Không service nào
  phụ thuộc \code{micrometer-registry-prometheus}, nên Actuator không
  có gì để lộ ra dưới tên đó. Thêm đúng một artifact ấy (BOM Boot quản
  lý phiên bản) làm dòng exposure hiện có thành thật.
  \item Độ trễ: không đổi --- span được xếp hàng và gửi bởi thread nền.
  Heap: bị chặn --- hàng đợi giữ 10.000 span và bỏ khi tràn kèm cảnh
  báo có đếm. Span: những cái sinh ra trong giờ đó mất; không có
  replay. Với reporter đồng bộ, \emph{độ trễ} sẽ đổi: mọi request sẽ
  chặn trên một lời gọi HTTP thất bại tới Zipkin cho tới khi timeout.
  \item Một phần trăm --- lấy mẫu head-based quyết định tại gốc trước
  khi biết kết quả, nên lỗi được giữ cùng tỷ lệ với thành công. Nâng
  lên 1.0 chỉ cho lỗi đòi lấy mẫu tail-based, thứ cần một OpenTelemetry
  Collector đệm trọn trace trước kho; đường trực tiếp Brave-tới-Zipkin
  không làm được.
  \item E-mail là giá trị tag không giới hạn; mỗi giá trị khác nhau tạo
  một chuỗi thời gian mới trong registry metric và một giá trị được
  đánh index mới trong kho trace. Bộ nhớ lớn theo số người dùng cho tới
  khi \code{/actuator/metrics} bò lê và pod bị \code{OOMKilled}. Quy
  tắc: tag mang \emph{mẫu và enum} (mẫu route, status, outcome); định
  danh đi vào annotation của span hay dòng log, những thứ không được
  đánh index theo giá trị.
\end{enumerate}

\section*{Chương 38 --- Đào sâu: Trình duyệt gặp gateway}
\begin{enumerate}
  \item Khóa \code{auth-storage} trong \code{localStorage}, do
  \code{persist} của Zustand ghi. Một dòng:
  \code{JSON.parse(localStorage.\allowbreak getItem('auth-storage')).\allowbreak
  state.refreshToken}.
  Có nó, script (hay bất kỳ ai nó gửi tới) có thể gọi
  \code{/api/auth/refresh} suốt bảy ngày và đúc access token mới tùy
  ý, từ bất kỳ máy nào.
  \item Một 401 từ \code{/refresh} đi qua \code{apiClient} sẽ vào
  interceptor phản hồi; \code{\_retry} chưa đặt trên request mới đó,
  nên nó sẽ đọc refresh token và POST \code{/refresh} lần nữa qua
  \code{apiClient}; 401 đó lại vào interceptor --- một chuỗi lời gọi
  refresh không giới hạn, mỗi cái sinh ra cái tiếp, cho tới khi đóng
  tab. \code{axios} trần không có interceptor, nên 401 của nó tới
  \code{catch} và đăng xuất.
  \item Cả bốn nhận 401 và mỗi cái tự đặt \code{\_retry}; bốn POST tới
  \code{/refresh} đua nhau. Cái đầu thành công và xoay token; cái thứ
  hai đưa ra token cũ giờ đã vô hiệu và nhận 401 từ auth-service; nó
  rơi vào \code{catch} --- bước 4, \code{logout()} và chuyển hướng
  cứng --- trong khi lần replay của request đầu thậm chí có thể thành
  công. Single-flight: một promise, khởi động bởi 401 đầu tiên, các
  cái còn lại await nó; một POST, một lần xoay, bốn lần replay cùng
  token mới.
  \item (a) Một endpoint có lỗ hổng hay bị lãng quên ở nơi khác dưới
  \code{/api} không bao giờ nhận cookie, nên lỗi CSRF hay rò rỉ log ở
  đó không chạm được nó; (b) cookie không đính kèm các lời gọi API
  thường, nên nó không bao giờ xuất hiện trong access log của gateway,
  proxy hay lưu lượng trình duyệt cho \code{/api/courses}. Trình duyệt
  so \code{Path} với URL \emph{nó} yêu cầu, tức
  \code{/api/auth/refresh} của gateway; \code{/refresh} đã bỏ tiền tố
  của auth-service không bao giờ là đường dẫn trình duyệt thấy.
  \item Không mode nào: proxy dev của Vite và Traefik trong cụm đều
  giữ trang và API trên một origin, nên không request nào là
  cross-origin và không preflight nào xảy ra. Phục vụ bundle đã build
  từ một origin khác --- một CDN, hay một static host ở cổng khác không
  qua proxy --- sẽ làm mọi lời gọi API thành cross-origin và cấu hình
  đó thành quyết định.
  \item CORS được trình duyệt cưỡng chế sau khi phản hồi tới: gateway
  đã xử lý request và trả 200, rồi trình duyệt từ chối trao body cho
  script vì phản hồi thiếu \code{Access-Control-Allow-Origin} chấp
  nhận được. So sánh: header \code{Origin} của request với danh sách
  allowed-origins (chuỗi chính xác, cả scheme và cổng);
  \code{Access-Control-Allow-Credentials} với việc request có
  credentialed hay không; và cho preflight,
  \code{Access-Control-Allow-Headers} với các header request thực sự
  gửi (\code{Authorization}, \code{Content-Type}).
\end{enumerate}
